\section{厚生比較の符号certificate}
\label{app:welfare}

本Appendixでは、Section~\ref{sec:coalition-stability} で用いるcanonical welfare comparisonについて解析的なsign certificateを与える。これらはStage-8 modelにおけるfrozen relationshipである。以下の計算はnumerical samplingに依存するのではなく、domain全体での符号の根拠を明示する。

次を再掲する。
\begin{align*}
W^{IS}&=K_I+3P,
& W_M^{SU,N}&=K_M+2A+C,\\
W_O^{SU,N}&=K_O+2B+D,
& W^{SW}&=K_W+H+2S.
\end{align*}
修正済みのinternational-standardization expressionは
\[
W^{IS}=\frac{3(5-2v)}{32(1-v)^2}.
\]
である。以下のすべての分母はcanonical domain上で厳格に正である。

\subsection{SU加盟国と各国別標準の比較}

scaled welfare gapを
\begin{align}
N_{MS}(c,v)
&\equiv 32(2-3v)^2\left(W_M^{SU,N}-W^{SW}\right)\notag\\
&=-216c^2v^2+264c^2v-60c^2
+108cv^2-144cv+56c
-99v^2+84v.
\label{eq:appendix-n-ms}
\end{align}
と定義する。$c$ に関するcurvatureは
\[
\frac{\partial^2N_{MS}}{\partial c^2}
=-24(18v^2-22v+5)<0,
\]
である。$18v^2-22v+5$ は $0<v<1/4$ で減少し、$v=1/4$ で $5/8$ だからである。したがって、許容される $c$ 区間のclosure上のminimumはendpointで達成される。$c=0$ では
\[
N_{MS}(0,v)=3v(28-33v)>0.
\]
である。$c_{\max}(v)=(1-3v)/[3(1-v)]$ では
\[
N_{MS}(c_{\max},v)
=-\frac{p_{MS}(v)}{3(1-v)^2},
\]
であり、ここで
\[
p_{MS}(v)=621v^4-1206v^3+729v^2-92v-36.
\]
である。$0\le v\le1/4$ では $v^4\le v^2/16$ であり、項 $-1206v^3$ は非正なので、
\[
p_{MS}(v)
\le \frac{12285}{16}v^2-92v-36.
\]
を得る。右辺の二次式はconvexで、$[0,1/4]$ の両endpointで負である。$v=1/4$ では $-2819/256$ である。したがって $p_{MS}(v)<0$ であり、
\begin{equation}
\boxed{W_M^{SU,N}>W^{SW}.}
\label{eq:appendix-su-member-gt-sw}
\end{equation}
を得る。

\subsection{国際標準化と各国別標準の比較}

同様に、
\begin{align}
N_{IS}(c,v)
&\equiv 32(1-v)^2\left(W^{IS}-W^{SW}\right)\notag\\
&=-28c^2v^2+56c^2v-28c^2
+20cv^2-40cv+20c
-15v^2+24v.
\label{eq:appendix-n-is}
\end{align}
とする。
\[
\frac{\partial^2N_{IS}}{\partial c^2}
=-56(1-v)^2<0,
\]
なので、その $c$ に関するminimumもendpointにある。endpoint valueは
\[
N_{IS}(0,v)=3v(8-5v)>0
\]
および
\[
N_{IS}(c_{\max},v)
=\frac{32+144v-207v^2}{9}>0.
\]
である。最後の不等式は $v<1/4$ から直ちに従う。$207v^2<207/16<32$ であり、$144v>0$ だからである。したがって、
\begin{equation}
\boxed{W^{IS}>W^{SW}.}
\label{eq:appendix-is-gt-sw}
\end{equation}
である。

\subsection{国際標準化とSU域外国の比較}

最後に、
\begin{align}
N_{IO}(c,v)
&\equiv 32(1-v)^2(2-3v)^2
\left(W^{IS}-W_O^{SU,N}\right).
\label{eq:appendix-n-io-definition}
\end{align}
と定義する。exact polynomialは
\begin{align}
N_{IO}(c,v)
={}&-144c^2v^4+576c^2v^3-912c^2v^2+672c^2v-192c^2\notag\\
&+288cv^4-912cv^3+1072cv^2-560cv+112c\notag\\
&-252v^4+666v^3-537v^2+132v.
\label{eq:appendix-n-io}
\end{align}
である。そのcurvatureは
\[
\frac{\partial^2N_{IO}}{\partial c^2}
=-96(1-v)^2(3v^2-6v+4)<0.
\]
である。$c=0$ では
\[
N_{IO}(0,v)
=-3v\,p_0(v),
\qquad
p_0(v)=84v^3-222v^2+179v-44.
\]
となる。導関数 $p_0'(v)=252v^2-444v+179$ は $[0,1/4]$ 上で減少し、$v=1/4$ でも $335/4$ と正である。したがって $p_0$ は増加するが、$p_0(1/4)=-189/16<0$ であるから、$N_{IO}(0,v)>0$ である。

$c$ の上限境界では、
\[
N_{IO}(c_{\max},v)
=-\frac{p_1(v)}{3},
\]
であり、
\[
p_1(v)=324v^4-990v^3+843v^2-92v-48.
\]
である。$v^4\le v^2/16$ を用い、非正の三次項を落とすと、
\[
p_1(v)
\le \frac{3453}{4}v^2-92v-48.
\]
となる。このbounding quadraticはconvexで、$[0,1/4]$ の両endpointで負である。$v=1/4$ では $-1091/64$ である。したがって $p_1(v)<0$ であり、$N_{IO}$ の上限境界値は正である。$c$ に関するconcavityから、
\begin{equation}
\boxed{W^{IS}>W_O^{SU,N}.}
\label{eq:appendix-is-gt-su-outsider}
\end{equation}
を得る。

\subsection{迂回後の域外国利得差}
\label{app:j-positive}

Section~\ref{sec:coalition-stability} では
\[
J=K_I+P-K_O-D,
\]
と定義したので、$W^{IS}-W_O^{SU,O}=J+F$ である。迂回後の域外国の厳格な順位付けを確立するには $J>0$ を示せば十分である。$J$ に正の分母を掛けると、符号同値な分子
\begin{align}
N_J(c,v)
={}&-48(1-v)^2c^2
+(48v^3-80v^2+16v+16)c\notag\\
&+v(68-269v+318v^2-108v^3).
\label{eq:appendix-j-numerator}
\end{align}
を得る。$v$ を固定すると、
\[
\frac{\partial^2N_J}{\partial c^2}=-96(1-v)^2<0.
\]
なので $N_J$ は $c$ についてconcaveである。したがって、許容される $c$ 区間のclosure上のminimumはendpointにある。$c=0$ と $c=c_{\max}(v)$ では、
\begin{align}
N_J(0,v)
&=v(68-269v+318v^2-108v^3),\notag\\
N_J(c_{\max},v)
&=\frac v3(284-1095v+1098v^2-324v^3).
\label{eq:appendix-j-endpoints}
\end{align}
である。括弧内の両polynomialは $0<v<1/4$ で減少し、$v=1/4$ でもそれぞれ $303/16$ と $1181/16$ で正である。したがって $N_J>0$ であり、ゆえに
\begin{equation}
\boxed{J>0.}
\label{eq:appendix-j-positive}
\end{equation}
である。$F>0$ なので、域外国企業のみの継続状態が関連する場合、canonical domain全体で $W^{IS}>W_O^{SU,O}$ が従う。

Equations~\eqref{eq:appendix-su-member-gt-sw}--\eqref{eq:appendix-j-positive} は、strict-blocking analysisで使用した符号比較そのものである。symbolic verification scriptは、manuscript textとは独立に、scaled-polynomial representation、curvature expression、endpoint identityを確認する。
